\section{动态规划进阶：参数、调用与改造}

本节紧跟前面的基础模板。先看每个接口传什么，再看哪些不变量不能动，最后看如何替换维护信息或转移。

\subsection{复杂数位动态规划与自动机数位动态规划}

数位 DP 处理 \passthrough{\lstinline!0..N!} 中满足数字条件的数。状态通常是 \passthrough{\lstinline!(pos,tight,started,extra)!}；上下界用 \passthrough{\lstinline!count(R)-count(L-1)!}。若有``不能出现某模式/相邻数字限制''，把 \passthrough{\lstinline!extra!} 换成自动机节点。

\begin{lstlisting}[language={C++}]
long long memo[20][2][2][STATE]; // -1=未计算；末维 st 为题目/自动机状态 [0,STATE)。
// 接口：dfs(pos, tight, started, st)：从当前节点/位置继续深度优先处理，并写入本节状态数组；返回类型 long long。
// 参数：当前 0-based/DP 位置 pos；前缀是否仍贴住上界 tight；是否已经放置非前导零数字 started；题目定义或自动机节点状态 st，范围 [0,STATE)。
long long dfs(int pos, bool tight, bool started, int st) {
    // pos：正在处理第 pos 位；tight：前缀是否仍等于上界；
    // started：是否已经出现过非前导零；st：额外限制的自动机状态。
    if (pos == len) return valid(st, started);
    // 只有 tight=false 的状态可以复用；tight=true 还依赖上界前缀。
    if (!tight && memo[pos][started][0][st] != -1)
        return memo[pos][started][0][st];
    // 变量：res=准备返回的结果容器/结果值 res。
    long long res = 0;
    // 变量：up=父侧贡献/倍增祖先数组 up。
    int up = tight ? digit[pos] : 9;
    // 变量：d=当前距离、差值或覆盖增量 d。
    for (int d = 0; d <= up; ++d) {
        // 前导零通常不应触发“数字出现/相邻数字”等条件。
        // 变量：ns=转移后的 started 状态 ns。
        bool ns = started || d != 0;
        // 尚未开始时是否把 0 送入自动机，必须按题意决定；这里留在起点 0。
        // 变量：nst=转移后的自动机/DP 状态 nst。
        int nst = ns ? go[st][d] : 0;
        if (ns && nst == forbidden) continue;
        // 注意必须与真实上界位 digit[pos] 比较，不能写成 d==up。
        res += dfs(pos + 1, tight && d == digit[pos], ns, nst);
    }
    if (!tight) memo[pos][started][0][st] = res;
    return res;
}
\end{lstlisting}

\subsubsection{测试样例：数位 DP 的前导零和上下界差分}

这里固定 \passthrough{\lstinline!valid!} 为``数字十进制表示中不出现数字 \passthrough{\lstinline!4!}''，并把 \passthrough{\lstinline!0!} 计入合法数。每组输入 \passthrough{\lstinline!L R!}，输出 \passthrough{\lstinline![L,R]!} 中合法整数个数。

\begin{lstlisting}
输入
4
0 0
0 4
0 10
0 99

输出
1
4
10
81
\end{lstlisting}

\passthrough{\lstinline![0,4]!} 的合法数是 \passthrough{\lstinline!0,1,2,3!}；\passthrough{\lstinline![0,99]!} 的答案是 \passthrough{\lstinline!9*9=81!}。如果把前导零当成真实数字 \passthrough{\lstinline!4!} 的相邻字符，或者把 \passthrough{\lstinline!tight \&\& d == digit[pos]!} 错写成 \passthrough{\lstinline!d == up!}，这些边界会出错。

这里 \passthrough{\lstinline!up!} 只表示当前位允许枚举到哪里；tight 更新必须比较真实上界位 \passthrough{\lstinline!digit[pos]!}。练习：\href{https://www.luogu.com.cn/problem/P2602}{P2602 数字计数}；若题目问每位出现次数，把返回值从一个数改为``方案数 + 各位贡献''的结构体。


\subsection{所有子集和动态规划、轮廓线动态规划与状态压缩}

SOS DP 用来计算所有子集/超集的和：

\begin{lstlisting}[language={C++}]
// 变量：bit=当前处理的二进制位 bit。
for (int bit = 0; bit < K; ++bit)
    // 变量：mask=当前子集/博弈状态的二进制掩码 mask。
    for (int mask = 0; mask < (1 << K); ++mask)
        if (mask >> bit & 1)
            // 把“去掉 bit 后的子集”贡献给当前 mask。
            // 处理完这一层后，f[mask] 已包含所有只在低 bit 位不同的子集。
            f[mask] += f[mask ^ (1 << bit)];
\end{lstlisting}

\subsubsection{测试样例：SOS 动态规划的子集和方向}

下面约定 \passthrough{\lstinline!K=2!}，初始 \passthrough{\lstinline!f[mask]!} 是恰好等于 \passthrough{\lstinline!mask!} 的权值，输出所有子集和 \passthrough{\lstinline!sum[sub]!}。

\begin{lstlisting}
输入
2
1
5 7
2
1 2 3 4

输出
5 12
1 3 4 10
\end{lstlisting}

第二组中 \passthrough{\lstinline!sum[3]!} 必须是 \passthrough{\lstinline!1+2+3+4=10!}；如果把 \passthrough{\lstinline!mask \^ (1<<bit)!} 写成加上不存在的超集，或循环顺序和转移方向不匹配，会在 \passthrough{\lstinline!sum[1]!}、\passthrough{\lstinline!sum[2]!} 同时出错。

若网格宽度 \passthrough{\lstinline!m<=10/12!}，每一列的填法只影响相邻列，维护 \passthrough{\lstinline!dp[col][mask]!}；这就是轮廓线 DP。先预处理从 \passthrough{\lstinline!mask!} 到 \passthrough{\lstinline!next\_mask!} 的合法转移，再做列 DP。棋盘放置题可练\href{https://www.luogu.com.cn/problem/P2704}{P2704 炮兵阵地}。状态中相邻两行都可能影响当前行时，不能只保留一行。


\subsection{单调队列、斜率优化与李超树优化动态规划}

若转移有 \passthrough{\lstinline!dp[i] = min(dp[j] + w(j,i))!}，且 \passthrough{\lstinline!j!} 的优劣具有区间单调性，先尝试单调队列；若整理后是直线最值，才上凸包/李超树。

\begin{lstlisting}[language={C++}]
// 变量：q=当前 BFS/单调算法使用的队列 q。
deque<int> q;
q.push_back(0); // 0 是已计算的初始候选；调用前先设置 dp[0]。
// 变量：i=当前循环下标 i。
for (int i = 1; i < n; ++i) {
    // 过期候选先删除；窗口条件由题目决定。
    while (!q.empty() && q.front() < i - window) q.pop_front();

    // 队首是当前窗口中最优候选，value(j,i) 是题目转移代价。
    dp[i] = value(q.front(), i);

    // 新候选 i 若在未来不可能优于队尾，就淘汰队尾。
    // next_x 和 value 的单调性必须由题目转移证明，不能盲套。
    while (!q.empty() && value(i, next_x(q.back()))
                          <= value(q.back(), next_x(q.back()))) q.pop_back();
    q.push_back(i);
}
\end{lstlisting}

\begin{lstlisting}[language={C++}]
// 把 dp[i] 整理成“查询某条候选直线在 x[i] 处的最值”时使用。
struct Line {
    long long k, b; // 候选直线 y=k*x+b 的斜率与截距。
    // 接口：value(x)：返回直线 k*x+b 在横坐标 x 处的值；返回类型 long long。
    // 参数：待处理值或点/状态 x。
    long long value(long long x) const { return k * x + b; }
};

// 只有在斜率按单调顺序加入、查询 x 也按单调顺序时，
// 才能用普通 deque；任意插入/任意查询请改用前面的李超线段树。
// 变量：hull=按有效顺序维护的凸包候选队列 hull。
deque<Line> hull;
// 接口：useless(a, b, c)：判断中间候选直线 b 是否可从凸包队列删除；返回 true 表示本节条件成立/操作成功。
// 参数：按斜率顺序相邻的候选直线 a；按斜率顺序相邻的候选直线 b；按斜率顺序相邻的候选直线 c。
bool useless(const Line& a, const Line& b, const Line& c) {
    // 交点顺序不再递增时，中间直线永远不会成为最优。
    // 用 __int128 防止交叉相乘时溢出。
    return (__int128)(b.b - a.b) * (b.k - c.k)
         >= (__int128)(c.b - b.b) * (a.k - b.k);
}

// 接口：add_line(line)：把一条候选直线插入当前李超树/单调凸包并删除失效候选；会修改当前结构；多测时重新构造或显式清空成员。
// 参数：待插入候选直线 line。
void add_line(Line line) {
    // 斜率顺序必须与 useless 的不等式方向一致；换成求最大值时通常要反号。
    while (hull.size() >= 2 &&
           useless(hull[hull.size() - 2], hull.back(), line))
        hull.pop_back();
    hull.push_back(line);
}
\end{lstlisting}

\subsubsection{调用测试：斜率优化只能在顺序条件成立时使用}

对下面三条直线做最小值查询：

\begin{lstlisting}
加入顺序：y=0，y=x，y=2x-1
查询顺序：x=0,1,3
应得到：0,0,1
\end{lstlisting}

再把查询顺序改成 \passthrough{\lstinline!x=3,0,1!}。如果模板仍然使用只支持单调查询的双端队列，必须重新检查题目是否真的满足单调性；不能把这组非单调调用当成普通凸包优化的合法输入。

斜率优化的关键不是代码，而是把转移整理为 \passthrough{\lstinline!y = kx+b!}，并证明斜率和查询点单调；若斜率不单调或要在线插入，直接使用前面的李超线段树。例题：\href{https://www.luogu.com.cn/problem/P5785}{P5785 任务安排}；先把平方项展开，再判断是否能维护直线。


\subsection{分治优化、四边形不等式、树形背包与 Steiner 树状压动态规划}

\subsubsection{调用测试：分治优化的决策范围必须覆盖真实最优点}

选一个可以暴力验证的代价：\passthrough{\lstinline!cost(j,i)=(prefix[i]-prefix[j])\^2!}，对 \passthrough{\lstinline!prefix=[0,1,3,6]!} 计算每一层的最优决策，并与 \passthrough{\lstinline!divide\_conquer\_dp!} 的结果逐项比较。重点不是某个固定答案，而是确认：

\begin{lstlisting}
测试一：optL=0,optR=i，结果应与完整枚举一致。
测试二：把 optL/optR 缩到不包含真实最优 j 的范围，程序必须被判定为调用错误，不能把错误范围当成算法结论。
\end{lstlisting}

没有单调性证明时，不能因为测试一通过就直接把完整枚举范围缩成 \passthrough{\lstinline!optL..optR!}。

分治优化要求最优决策点满足单调性：\passthrough{\lstinline!opt[l][r-1] <= opt[l][r]!}。模板递归 \passthrough{\lstinline!solve(L,R,optL,optR)!}，只在允许区间枚举决策点；没有证明时不能硬套。四边形不等式和决策单调性是证明工具，不是看到区间 DP 就默认成立。

\begin{lstlisting}[language={C++}]
// 典型形式：dp[layer][i] = min_j(dp[layer-1][j] + cost(j,i))。
// optL/optR 是当前区间内“最优决策 j”可能出现的范围。
// 接口：divide_conquer_dp(L, R, optL, optR)：计算 DP 下标闭区间 [L,R]，决策只搜索 [optL,optR]；无返回值；结果写入引用参数或当前结构状态。
// 参数：闭区间左端点 L；闭区间右端点 R；真实最优决策允许范围左端 optL；真实最优决策允许范围右端 optR。
void divide_conquer_dp(int L, int R, int optL, int optR) {
    if (L > R) return;
    // 变量：mid=当前区间中点 mid。
    int mid = (L + R) >> 1;
    // 变量：best=目前找到的最优值 best。
    long long best = INF;
    // 变量：best_j=当前 DP 状态取得最优值的决策下标 best_j。
    int best_j = -1;

    // 没有单调性证明时不能把枚举范围缩成 optL..optR；
    // 这份模板的正确性前提就是 opt[mid] 位于该范围内。
    // 变量：j=内层循环下标 j。
    for (int j = optL; j <= min(optR, mid); ++j) {
        // 变量：cand=当前枚举得到的候选答案 cand。
        long long cand = old_dp[j] + cost(j, mid);
        if (cand < best) best = cand, best_j = j;
    }
    new_dp[mid] = best;

    // 由 opt[l] <= opt[mid] <= opt[r]，左右区间可以分别缩小。
    divide_conquer_dp(L, mid - 1, optL, best_j);
    divide_conquer_dp(mid + 1, R, best_j, optR);
}
\end{lstlisting}

树形背包的状态是``子树里选了 \passthrough{\lstinline!j!} 个''，合并子树时做背包卷积；先处理小子树，循环边界必须使用当前已合并大小。Steiner 树状压 DP 适用于关键点数 \passthrough{\lstinline!k<=10/15!} 的``连通所有关键点最小代价''：先对每个 \passthrough{\lstinline!mask!} 做子集合并，再以 Dijkstra 在图上松弛。若关键点数超过二十，通常要重新分析题意。

\begin{lstlisting}[language={C++}]
// Steiner 树状态：dp[mask][v] = 连接 mask 中关键点，并以 v 作为当前末端的最小代价。
// 变量：mask=当前子集/博弈状态的二进制掩码 mask。
for (int mask = 1; mask < (1 << K); ++mask) {
    // 把 mask 拆成两个非空子集，在同一个 v 汇合。
    // 变量：sub=当前枚举的子掩码 sub。
    for (int sub = (mask - 1) & mask; sub; sub = (sub - 1) & mask) {
        // 变量：other=当前移动后的另一堆/另一状态 other。
        int other = mask ^ sub;
        if (sub > other) continue; // 对称拆分只做一次，常数更小。
        // 变量：v=当前相邻顶点/下一状态 v。
        for (int v = 0; v < n; ++v)
            dp[mask][v] = min(dp[mask][v],
                              dp[sub][v] + dp[other][v]);
    }

    // 合并后，末端 v 还可以沿图上的边移动；
    // 用 Dijkstra 在“固定 mask 的状态层”上继续松弛。
    dijkstra_on_state(mask);
}
\end{lstlisting}

动态 DP/矩阵线段树在本章 2.13 已给出调用思路；概率 DP 遇到自环时，把方程写成 \passthrough{\lstinline!f = p*f + rest!}，得到 \passthrough{\lstinline!f=rest/(1-p)!}；多个状态互相转移才需要高斯消元。不要在有自环时直接按拓扑序 DP。


\subsection{动态动态规划与矩阵线段树}

当树/序列上的一个点被修改后，答案仍然可以由局部状态转移合并，但普通 DP 不能重新算全局，就把每个位置的转移写成矩阵，线段树维护矩阵乘积。

\begin{lstlisting}[language={C++}]
struct Mat {
    long long a[2][2]{}; // a[i][j]：第 i 个状态对第 j 个状态的线性系数。
};
// 接口：mul(A, B)：返回两个转移矩阵的乘积；返回类型 Mat。
// 参数：左操作数/矩阵 A；右操作数/矩阵 B。
Mat mul(Mat A, Mat B) {
    Mat C{}; // 必须清零；矩阵乘法是累加，不是覆盖。
    // 变量：i=当前循环下标 i。
    for (int i = 0; i < 2; ++i)
        // 变量：k=当前排名、选取数量或循环层数 k。
        for (int k = 0; k < 2; ++k)
            // 变量：j=内层循环下标 j。
            for (int j = 0; j < 2; ++j)
                // C = A * B：先应用 B，再应用 A（按列向量约定）。
                C.a[i][j] = (C.a[i][j] + A.a[i][k] * B.a[k][j]) % MOD;
    return C;
}

// 接口：identity()：返回当前状态维数的单位转移矩阵；返回类型 Mat。
Mat identity() {
    // 变量：I=单位矩阵/恒等转移 I。
    Mat I{};
    I.a[0][0] = I.a[1][1] = 1;
    return I;
}

// 矩阵线段树的核心合并：左区间的转移先于右区间的转移。
// 如果状态向量按行向量存储，合并顺序会反过来，必须从建模时保持一致。
// 接口：merge_segment(left, right)：按从左到右的应用顺序合并两个区间转移矩阵；返回类型 Mat。
// 参数：左区间转移矩阵 left；右区间转移矩阵 right。
Mat merge_segment(Mat left, Mat right) {
    return mul(right, left);
}
\end{lstlisting}

现场映射方法：先在纸上写出``左边 DP 状态向右移动一格如何变成新状态''，把这一步写成矩阵；叶子存单点矩阵，线段树父节点存 \passthrough{\lstinline!left * right!}。矩阵乘法不可交换，区间方向不能写反。它是 C 级内容，不建议只看一遍标题就现场发明。


\subsection{概率动态规划的自环方程}

有自环的状态不能按拓扑序处理。例如从状态 \passthrough{\lstinline!i!} 以概率 \passthrough{\lstinline!p!} 留在 \passthrough{\lstinline!i!}，以概率 \passthrough{\lstinline!q!} 转到已知状态，期望满足：

\[E_i=pE_i+qE_j+c\quad\Longrightarrow\quad E_i=\frac{qE_j+c}{1-p}.\]

多个状态互相转移时，把每个状态写成一行：

\[E_i-\sum_j p_{ij}E_j=c_i,\]

然后用高斯消元。\passthrough{\lstinline!1-p!} 很小时误差会放大，题目若要求模意义下概率，要先确认所有分母可逆；题目若允许无限期望，要识别并输出题目规定的无穷值。

\subsubsection{测试样例：概率自环的除法和无限期望}

把``每次转移产生的固定代价''也写进方程，测试不能只测没有自环的普通转移：

\begin{lstlisting}
调用 1：E = 0.5 E + 1
期望：E = 2

调用 2：E = 0.5 E + 0.5 * 4 + 1
期望：E = 6

调用 3：E = 1.0 E + 1
期望：不存在有限值；不能直接计算 (1-p) 的逆。

调用 4：
  E1 = 0.5 E1 + 0.5 E2 + 1
  E2 = 0.25 E1 + 0.75 E2 + 2
期望：先移项得到
  0.5 E1 - 0.5 E2 = 1
 -0.25 E1 + 0.25 E2 = 2，
  两行线性相关但常数不相容，因此该系统无有限解。
\end{lstlisting}

前三组分别抓``漏除以 \passthrough{\lstinline!1-p!}''\,``把常数项放错侧''和``\passthrough{\lstinline!p=1!} 除零''。第四组抓高斯消元中的主元选择、无解判断与浮点误差。实际题目若保证期望存在，第四组应改成相容系统；测试代码本身仍应保留无解分支，避免模板在异常数据上产生假答案。
